\section{Additional Control Trials for Biotic Background Experiments}

To better characterize the mechanism behind fitness effects caused by the background strain, we performed additional specimen adaptation assays under biotic background conditions with diversity maintenance disabled.
This analysis allowed us to test whether action of the diversity maintenance mechanism, rather than direct interactions between the focal and background strains, caused the observed fitness effects.
Supplementary Figure \ref{fig:with_vs_without_diversity_maintenance} compares adaptation assay outcomes with and without diversity maintenance under both the prefatory and contemporary biotic background conditions.
Outcomes were generally similar, and we observed only one sign-change difference: one specimen outcome was beneficial under prefatory biotic background conditions without diversity maintenance, but deleterious with diversity maintenance.
Further, without diversity maintenance, we still observed four outcomes that were advantageous under biotic conditions but deleterious under abiotic conditions (Supplementary Figure \ref{fig:outcome_count_joint_distns:specimen_no_dm}).
So, biotic selective effects cannot be explained as an artifact of the diversity maintenance scheme.

We also conducted specimen adaptation assays with diversity maintenance disabled under the control focal strain biotic background.
In these experiments, we again found no evidence for impact from the diversity maintenance scheme on results
(Supplementary Figure \ref{fig:baseline_fitness_gain_or_loss}).
